\documentclass{article}
\usepackage{amsmath,amssymb,amsthm}
\usepackage{hyperref}
\newtheorem{theorem}{Theorem}[section]

% Compile with: pdflatex references; bibtex references; pdflatex references;
% pdflatex references. The matching database is references.bib.
\begin{document}
\tableofcontents
\section{Cross-references\label{sec:references}}

\LaTeX{} keeps references connected to their labels. This is
Section~\ref{sec:references}; \autoref{thm:identity} and
Equation~\eqref{eq:identity} below can be cited before their definitions.
Use \nameref{sec:references} for a heading name and
page~\pageref{eq:identity} for an object's page number.
Change a label in PasTeX to update the commands that refer to it.

\begin{theorem}[An elementary identity]\label{thm:identity}
For every real number $x$, the following identity holds.
\begin{equation}\label{eq:identity}
(x+1)^2=x^2+2x+1.
\end{equation}
\end{theorem}

\begin{equation}\tag{A}\label{eq:tagged}
\sum_{n=1}^{\infty}\frac{1}{n^2}=\frac{\pi^2}{6}.
\end{equation}
The manually tagged formula is Equation~\eqref{eq:tagged}.

\section{Bibliography citations}
The typesetting system is described in~\cite{knuth1984}, and
\LaTeX{} document preparation in~\cite{lamport1994}.
Several entries can be cited together: \cite{knuth1984,lamport1994}.
A page note can be added: \cite[Chapter~2]{lamport1994}.

\begin{description}
\item[Labels] Use \verb|\ref| for sections and \verb|\eqref| for equations.
\item[Citations] Use \verb|\cite| to refer to entries in the bibliography.
\end{description}

\section{Try the reference tools}
Click a reference to choose a different target, and edit a label beside its
object to update its uses. Ctrl+hover underlines actionable entries; Ctrl+click
follows the target. Normal clicks remain editing actions. Compare a numbered
reference, \ref{sec:references}, with its named form, \autoref{sec:references}.

Click a bibliography citation to edit its keys or page note. References,
Bibliography selects the database, backend and installed styles. This file uses
traditional BibTeX and the plain style; a saved copy can be used to try another
style or biblatex. Refresh the preview after changing the setup. The contents
and bibliography blocks are read only, but their displayed text can be copied;
Ctrl+click follows their links. Edit the source heading or database entry to
change generated text.

\bibliographystyle{plain}
\bibliography{references}
\end{document}
